% Research poster. Copy this folder to <paper project>/poster/ and fill it in.
% The poster is almost only figures. Write the message as a sentence or two, split them into
% parts, and give each part a card with the figure for it: the headings, read in order, make the
% sentences.
% Compile with pdflatex (Overleaf's default): latexmk -pdf poster.tex
\documentclass[final]{beamer}
% Paper size in cm: 48 x 36 in landscape is width=121.92,height=91.44; A0 portrait is
% width=84.1,height=118.9. scale sets the text size: body text is 24.9pt x scale (36pt here).
\usepackage[size=custom,width=121.92,height=91.44,scale=1.45]{beamerposter}
% The paper's main method colour, the colour its figures give it (default teal):
% \definecolor{posteraccent}{HTML}{007E80}
\usetheme{poster}
% Figures: \fig{../figs/name} enlarges a paper figure, labels included, to fill its card;
% \fig[.6]{...} fills at most .6 of the card's height, for a card with more in it. To redraw a
% figure at poster size instead, load the paper's figure style, then enlarge its fonts
% (\renewcommand\labeltype{\small}, \renewcommand\ticktype{\footnotesize}) and marks:
% \input{../figs/src/figstyle}

\title{Title of the Paper:}
\subtitle{Second Line of the Title}
\author{First Author, Second Author}
\institute{Department, The Ohio State University}
\posterlogo{osu-block-o}
% The conference's logo, a file beside poster.tex:
% \conferencelogo{conference-logo}

\begin{document}
\begin{frame}[t]

% Rows: the heights (.3, .35, .35) add up to 1; in each row, the cards' columns add up to 24
\begin{cardrow}{.3}
  \begin{herocard}{24}
    \begin{minipage}[c]{.36\linewidth}
      \placeholder{16cm}{Teaser figure}
    \end{minipage}\hfill
    \begin{minipage}[c]{.5\linewidth}
      \claim{The main claim of the paper, in one line.}
      \statement{The setting, with its \textbf{key terms} in bold.\\
        What the paper finds about \textbf{the first thing}.\\
        What it finds about \textbf{the second thing}.}
    \end{minipage}\hfill
    \begin{minipage}[c]{.1\linewidth}
      \qrlink{https://example.com}{Paper}
    \end{minipage}
  \end{herocard}
\end{cardrow}

\begin{cardrow}{.35}
  \begin{card}{8}{The first part of the sentence,}
    \placeholder{15cm}{Figure}
  \end{card}
  \begin{card}{10}{its second part,}
    \begin{outlinebox}{A definition}
      \[ s(x) = \frac{\lVert h(x) - \mu_- \rVert_2^2}{\lVert h(x) - \mu_+ \rVert_2^2} \]
    \end{outlinebox}
    \placeholder{8cm}{Figure}
  \end{card}
  \begin{card}{6}{and its end.}
    \placeholder{15cm}{Figure}
  \end{card}
\end{cardrow}

\begin{cardrow}{.35}
  \begin{card}{12}{A second sentence starts here,}
    \begin{theorem}[Name]
      The main result, stated as in the paper, for an input $x$.
    \end{theorem}
    \placeholder{9cm}{Figure}
  \end{card}
  \begin{card}{12}{and ends with the result!}
    \placeholder{14cm}{Figure}
  \end{card}
\end{cardrow}

\end{frame}
\end{document}
